\chapter{Background}
\label{ch:background}

\section{HPCAgent-Bench}
% Kernel = NumPy reference (oracle) + input generator + YAML manifest (presets S/M/L/XL,
% graded outputs). Elementwise grading |a - e| <= atol + rtol |e|. Agent campaigns
% (llr40v10 arms). Translator (numpyto_c) produces the C/C++/Fortran/numba lowerings.

\section{The Berkeley Dwarfs}
% Asanovic et al. 2006: 13 computational patterns; used to place the extracted kernels
% and to show what LLR-40 (loop-level, mostly structured) does not cover.

\section{Compiler Auto-Vectorization and Loop-Level Reasoning}
% TSVC (Callahan et al.), Maleki et al., TSVC-2. Loop transformations used as LLR-40
% classes: interchange, distribution, fusion, reduction, wavefront, scalar expansion, ...
% Cost models (Pohl et al.), limits of TSVC as an ideal case (Siso et al.).

\section{Floating-Point Semantics and Compiler Flags}
% -ffp-contract=fast (FMA), -fassociative-math / reassociation, why -ffast-math breaks an
% elementwise oracle. gfortran vs gcc reassociation difference.

\section{Code Representations}
% numba (JIT, parallel=True, no fastmath), C, C++, Fortran; translator-generated vs
% hand-written; DaCe and the SDFG as a further generator.

\section{Iterative Solvers and Preconditioning}
% Krylov methods (CG, GMRES), Gauss-Seidel/SOR, ILU/IC, sparse triangular solve and level
% scheduling, multigrid, Newton-Krylov/JFNK, time integrators.

\section{Agent Skills}
% What a skill is (SKILL.md, progressive disclosure, reference files), why it is the
% vehicle for domain knowledge an agent otherwise lacks.
